\ifdefined\gkver\else\def\gkver{student}\fi
\documentclass[\gkver]{gaokaozhenti}
\begin{document}

\chapter{2007年全国理综Ⅰ}
%% paperType: 理综物理部分

\begin{enumerate}
\item
%% number: 14
%% paperName: 2007年全国理综Ⅰ第 14 题 6 分
%% typeId: 1
%% type: 单选题
%% chapter: 曲线运动
%% point: 万有引力定律
%% method: 无
%% score: 6
%% degree: 650
%% duplicateId: 0
%% body:
据报道，最近在太阳系外发现了首颗“宜居”行星，其质量约为地球质量的 6.4 倍，一个在地球表面重量为 $600\UN$ 的人在这个行星表面的重量将变为 $960\UN$。由此可推知，该行星的半径与地球半径之比约为\xzanswer{B}
\fourchoices[answer=B]
{0.5}
{2}
{3.2}
{4}
%% answer: B
%% memo:
\memoanswer{
由题意可以得到 $g'=1.6g$；由黄金代换 $GM=gR^{2}$ 可以得到 $\frac{M'R^{2}}{MR'^{2}}=\frac{g'}{g}$，解得 $R'=2R$。
}

\item
%% number: 15
%% paperName: 2007年全国理综Ⅰ第 15 题 6 分
%% typeId: 1
%% type: 单选题
%% chapter: 机械波
%% point: 波的图象
%% method: 无
%% score: 6
%% degree: 650
%% duplicateId: 0
%% body:
一列简谐横波沿 $x$ 轴负方向传播，波速为 $v=4\Ums$。已知坐标原点（$x=0$）处质点的振动图像如图 a 所示，在下列 4 幅图中能够正确表示 $t=0.15\Us$ 时波形的图是\xzanswer{A}

\onepicture[num=a, w=6.0cm]{figs/15.png}

\fourchoices[answer=A, ispicture=true, h=2.0cm]
{figs/15a.png}
{figs/15b.png}
{figs/15c.png}
{figs/15d.png}
%% answer: A
%% memo:
\memoanswer{
由振动图像得到原点处的质点在 $y$ 正半轴向下运动，由于向负 $x$ 轴传播，所以只有 A 选项正确。
}

\item
%% number: 16
%% paperName: 2007年全国理综Ⅰ第 16 题 6 分
%% typeId: 2
%% type: 多选题
%% chapter: 气体性质
%% point: 热力学第一定律
%% method: 无
%% score: 6
%% degree: 650
%% duplicateId: 0
%% body:
如图所示，质量为 $m$ 的活塞将一定质量的气体封闭在气缸内，活塞与气缸之间无摩擦。a 态是气缸放在冰水混合物中气体达到的平衡状态，b 态是气缸从容器中移出后，在室温（$27\UduC$）中达到的平衡状态。气体从 a 态变化到 b 态的过程中大气压强保持不变。若忽略气体分子之间的势能，下列说法正确的是\xzanswer{AC}

\onepicture[w=4.5cm]{figs/16.png}

\fourchoices[answer=AC]
{与 b 态相比，a 态的气体分子在单位时间内撞击活塞的个数较多}
{与 a 态相比，b 态的气体分子在单位时间内对活塞的冲量较大}
{在相同时间内，a、b 两态的气体分子对活塞的冲量相等}
{从 a 态到 b 态，气体的内能增加，外界对气体做功，气体对外界释放了热量}
%% answer: AC
%% memo:
\memoanswer{
由于两种状态下压强相等，所以在单位时间单位面积里气体分子对活塞的总冲量肯定相等；由于 b 状态的温度比 a 状态的温度要高，所以分子的平均动量增大，因为总冲量保持不变，所以 b 状态单位时间内冲到活塞的分子数肯定比 a 状态要少。
}

\item
%% number: 17
%% paperName: 2007年全国理综Ⅰ第 17 题 6 分
%% typeId: 1
%% type: 单选题
%% chapter: 光学
%% point: 光的折射
%% method: 无
%% score: 6
%% degree: 650
%% duplicateId: 0
%% body:
在桌面上有一倒立的玻璃圆锥，其顶点恰好与桌面接触，圆锥的轴（图中虚线）与桌面垂直，过轴线的截面为等边三角形，如图所示。有一半径为 $r$ 的圆柱形平行光束垂直入射到圆锥的底面上，光束的中心轴与圆锥的轴重合。已知玻璃的折射率为 1.5，则光束在桌面上形成的光斑半径为\xzanswer{C}

\onepicture[w=5.0cm]{figs/17.png}

\fourchoices[answer=C]
{$r$}
{1.5$r$}
{2$r$}
{2.5$r$}
%% answer: C
%% memo:
\memoanswer{
如图所示，光线射到 A 或 B 时，入射角大于临界角，发生全反射，而后由几何关系得到第二次到达界面的时候垂直打出。O 点为 $\triangle ABC$ 的重心，设 $EC=x$，则由几何关系得到：$\frac{x}{x+r}=\frac{2}{3}$，解得光斑半径 $x=2r$。

\onepicture[w=4.5cm]{figs/17_an.jpg}
}

\item
%% number: 18
%% paperName: 2007年全国理综Ⅰ第 18 题 6 分
%% typeId: 1
%% type: 单选题
%% chapter: 运动和力
%% point: 动量定理
%% method: 无
%% score: 6
%% degree: 650
%% duplicateId: 0
%% body:
如图所示，在倾角为 $30^{\circ}$ 的足够长的光滑斜面上有一质量为 $m$ 的物体，它受到沿斜面方向的力 $F$ 的作用。力 $F$ 可按图（a）、（b）、（c）、（d）所示的四种方式随时间变化（图中纵坐标是 $F$ 与 $mg$ 的比值，力沿斜面向上为正）。已知此物体在 $t=0$ 时速度为零，若用 $v_{1}$、$v_{2}$、$v_{3}$、$v_{4}$ 分别表示上述四种受力情况下物体在 3 秒末的速率，则这四个速率中最大的是\xzanswer{C}

\onepicture[w=4.5cm]{figs/18a.png}

\onepicture[w=13.5cm]{figs/18b.png}

\fourchoices[answer=C]
{$v_{1}$}
{$v_{2}$}
{$v_{3}$}
{$v_{4}$}
%% answer: C
%% memo:
\memoanswer{
选向下为正方向，由动量定理分别得到

对于图（a）：$1.5mg\times2+0.5mg\times1=mv_{1}$

对于图（b）：$0.5mg\times1+mg\times1+1.5mg\times1=mv_{2}$

对于图（c）：$mg\times1+1.5mg\times2=mv_{3}$

对于图（d）：$1.5mg\times2=mv_{4}$

综合四个选项得到 $v_{3}$ 最大。
}

\item
%% number: 19
%% paperName: 2007年全国理综Ⅰ第 19 题 6 分
%% typeId: 2
%% type: 多选题
%% chapter: 原子和原子核
%% point: 玻尔模型
%% method: 无
%% score: 6
%% degree: 650
%% duplicateId: 0
%% body:
用大量具有一定能力的电子轰击大量处于基态的氢原子，观测到了一定数目的光谱线。调高电子的能力在此进行观测，发现光谱线的数目比原来增加了 5 条。用 $\Delta n$ 表示两侧观测中最高激发态的量子数 $n$ 之差，$E$ 表示调高后电子的能量。根据氢原子的能级图可以判断，$\Delta n$ 和 $E$ 的可能值为\xzanswer{AD}

\onepicture[w=5.0cm]{figs/19.png}

\fourchoices[answer=AD]
{$\Delta n=1$，$13.22\UeV<E<13.32\UeV$}
{$\Delta n=2$，$13.22\UeV<E<13.32\UeV$}
{$\Delta n=1$，$12.75\UeV<E<13.06\UeV$}
{$\Delta n=2$，$12.75\UeV<E<13.06\UeV$}
%% answer: AD
%% memo:
\memoanswer{
存在两种可能，第一种 $n=2$ 到 $n=4$，由于是电子轰击，所以电子的能量必须满足 $13.6\UeV-0.85\UeV<E<13.6\UeV-0.54\UeV$，故 D 选项正确；第二种可能是 $n=5$ 到 $n=6$，电子能量必须满足 $13.6\UeV-0.38\UeV<E<13.6\UeV-0.28\UeV$，故 A 选项正确。
}

\item
%% number: 20
%% paperName: 2007年全国理综Ⅰ第 20 题 6 分
%% typeId: 1
%% type: 单选题
%% chapter: 电场
%% point: 等势面
%% method: 无
%% score: 6
%% degree: 450
%% duplicateId: 0
%% body:
abcd 是匀强电场中的四个点，它们正好是一个矩形的四个顶点。电场线与矩形所在的平面平行。已知 a 点的电势是 $20\UV$，b 点的电势是 $24\UV$，d 点的电势是 $4\UV$，如图。由此可知，c 点的电势为\xzanswer{B}

\onepicture[w=6.0cm]{figs/20.png}

\fourchoices[answer=B]
{$4\UV$}
{$8\UV$}
{$12\UV$}
{$24\UV$}
%% answer: B
%% memo:
\memoanswer{
运用一个结论：在匀强电场中，任意一族平行线上等距离的两点的电势差相等，所以 $U_{ab}=U_{cd}$，所以 c 点电势为 $8\UV$。
}

\item
%% number: 21
%% paperName: 2007年全国理综Ⅰ第 21 题 6 分
%% typeId: 1
%% type: 单选题
%% chapter: 电磁感应
%% point: 电磁感应中图象
%% method: 无
%% score: 6
%% degree: 450
%% duplicateId: 0
%% body:
如图所示，LOO$'$L$'$ 为一折线，它所形成的两个角 $\angle$LOO$'$ 和 $\angle$OO$'$L$'$ 均为 $45^{\circ}$。折线的右边有一匀强磁场，其方向垂直于纸面向里，一边长为 $l$ 的正方形导线框沿垂直于 OO$'$ 的方向以速度 $v$ 做匀速直线运动，在 $t=0$ 时刻恰好位于图中所示的位置。以逆时针方向为导线框中电流的正方向，在下面四幅图中能够正确表示电流—时间（$I-t$）关系的是（时间以 $l/v$ 为单位）\xzanswer{D}

\onepicture[w=4.5cm]{figs/21.png}

\fourchoices[answer=D, ispicture=true, h=2.0cm]
{figs/21a.png}
{figs/21b.png}
{figs/21c.png}
{figs/21d.png}
%% answer: D
%% memo:
\memoanswer{
由初始位置可得，切割的有效长度在逐渐变大，且为逆时针，所以 BD 中选一个，由于 BD 两项中第 2 秒是一样的，没有区别。在第 3 秒内，线框已经有部分出上面磁场，切割的有效长度在减少，且为顺时针方向，所以只有 D 选项是正确的。
}

\item
%% number: 22
%% paperName: 2007年全国理综Ⅰ第 22 题 11 分
%% typeId: 4
%% type: 实验题
%% chapter: 运动和力
%% point: 动量守恒定律
%% method: 无
%% score: 11
%% degree: 450
%% duplicateId: 0
%% body:
\begin{enumerate}
	\item
	用示波器观察频率为 $900\UHz$ 的正弦电压信号。把该信号接入示波器 Y 输入。

	\onepicture[align=right, w=5.0cm]{figs/22a.png}

	\begin{steps}
		\item
		当屏幕上出现如图 1 所示的波形时，应调节 \tkanswer[1.5cm]{} 钮。如果正弦波的正负半周均超出了屏幕的范围，应调节 \tkanswer[1.5cm]{} 钮或 \tkanswer[1.5cm]{} 钮，或这两个钮配合使用，以使正弦波的整个波形出现在屏幕内。
		\item
		如需要屏幕上正好出现一个完整的正弦波形，应将 \tkanswer[1.5cm]{} 钮置于 \tkanswer[1.5cm]{} 位置，然后调节 \tkanswer[1.5cm]{} 钮。
	\end{steps}

	\onepicture[num=1, w=4.5cm]{figs/22b.png}

	\item
	碰撞的恢复系数的定义为 $e=\frac{|v_{2}-v_{1}|}{v_{20}-v_{10}}$，其中 $v_{10}$ 和 $v_{20}$ 分别是碰撞前两物体的速度，$v_{1}$ 和 $v_{2}$ 分别是碰撞后物体的速度。弹性碰撞的恢复系数 $e=1$，非弹性碰撞的 $e<1$。某同学借用验证动量守恒定律的实验装置（如图所示）验证弹性碰撞的恢复系数是否为 1，实验中使用半径相等的钢质小球 1 和 2（它们之间的碰撞可近似视为弹性碰撞），且小球 1 的质量大于小球 2 的质量。

	实验步骤如下：

	安装好实验装置，做好测量前的准备，并记下重锤线所指的位置 O。

	第一步，不放小球 2，让小球 1 从斜槽上 A 点由静止滚下，并落在地面上。重复多次，用尽可能小的圆把小球的所有落点圈在里面，其圆心就是小球落点的平均位置。

	第二步，把小球 2 放在斜槽前端边缘处 C 点，让小球 1 从 A 点由静止滚下，使它们碰撞。重复多次，并使用与第一步同样的方法分别标出碰撞后小球落点的平均位置。

	第三步，用刻度尺分别测量三个落地点的平均位置离 O 点的距离，即线段 OM、OP、ON 的长度。上述实验中，
	\begin{steps}
		\item
		P 点是 \tkanswer[2.0cm]{} 平均位置，M 点是 \tkanswer[2.0cm]{} 平均位置，N 点是 \tkanswer[2.0cm]{} 平均位置。
		\item
		请写出本实验的原理 \tkanswer[4.0cm]{}，写出用测量量表示的恢复系数的表达式 \tkanswer[4.0cm]{}。
		\item
		三个落地点距 O 点的距离 OM、OP、ON 与实验所用的小球质量是否有关系？\tkanswer[5.0cm]{}。
	\end{steps}

	\onepicture[align=right, w=6.0cm]{figs/22c.png}
\end{enumerate}
%% answer: 
\jdanswer{
\begin{enumerate}
	\item
	（1）竖直位移或$\uparrow\downarrow$，衰减或衰减调节，Y 增益；（2）扫描范围，1 k 挡位，扫描微调。

	\item
	（1）P 点是在实验的第一步中小球 1 落点的平均位置，M 点是小球 1 与小球 2 碰后小球 1 落点的平均位置，N 点是小球 2 落点的平均位置；（2）小球从槽口 C 飞出后做平抛运动的时间相同，设为 $t$，则 $\mathrm{OP}=v_{10}t$，$\mathrm{OM}=v_{1}t$，$\mathrm{ON}=v_{2}t$，小球 2 碰撞前静止，即 $v_{20}=0$，故 $e=\frac{v_{2}-v_{1}}{v_{10}-v_{20}}=\frac{\mathrm{ON}-\mathrm{OM}}{\mathrm{OP}}$；（3）OP 与小球的质量无关，OM 和 ON 与小球的质量有关。

\end{enumerate}
}
%% memo:
\memoanswer{
本题原卷及材料中未提供详解，待补充。
}

\item
%% number: 23
%% paperName: 2007年全国理综Ⅰ第 23 题 15 分
%% typeId: 6
%% type: 计算题
%% chapter: 直线运动
%% point: 相遇问题
%% method: 无
%% score: 15
%% degree: 650
%% duplicateId: 0
%% body:
甲、乙两运动员在训练交接棒的过程中发现：甲经短距离加速后能保持 $9\Ums$ 的速度跑完全程；乙从起跑后到接棒前的运动是匀加速的。为了确定乙起跑的时机，需在接力区前适当的位置设置标记。在某次练习中，甲在接力区前 $s_{0}=13.5\Um$ 处作了标记，并以 $v=9\Ums$ 的速度跑到此标记时向乙发出起跑口令。乙在接力区的前端听到口令时起跑，并恰好在速度达到与甲相同时被甲追上，完成交接棒。已知接力区的长度为 $L=20\Um$。求：
\begin{enumerate}
	\item
	此次练习中乙在接棒前的加速度 $a$；
	\item
	在完成交接棒时乙离接力区末端的距离。
\end{enumerate}
%% answer: 
\jdanswer{
\begin{enumerate}
	\item
	$3\Umsq$
	\item
	$6.5\Um$
\end{enumerate}
}
%% memo:
\memoanswer{
（1）设经过时间 $t$，甲追上乙，则根据题意有 $vt-\frac{1}{2}vt=13.5$。将 $v=9$ 代入得到：$t=3\Us$，再由 $v=at$ 解得：$a=3\Umsq$。

（2）在追上乙的时候，乙走的距离为 $s$，则：$s=\frac{1}{2}at^{2}$。代入数据得到 $s=13.5\Um$，所以乙离接力区末端的距离为 $\Delta s=20\Um-13.5\Um=6.5\Um$。
}

\item
%% number: 24
%% paperName: 2007年全国理综Ⅰ第 24 题 16 分
%% typeId: 6
%% type: 计算题
%% chapter: 运动和力
%% point: 动量守恒定律
%% method: 无
%% score: 16
%% degree: 450
%% duplicateId: 0
%% body:
如图所示，质量为 $m$ 的由绝缘材料制成的球与质量为 $M=19m$ 的金属球并排悬挂。现将绝缘球拉至与竖直方向成 $\theta=60^{\circ}$ 的位置自由释放，下摆后在最低点与金属球发生弹性碰撞。在平衡位置附近存在垂直于纸面的磁场。已知由于磁场的阻尼作用，金属球将于再次碰撞前停在最低点处。求经过几次碰撞后绝缘球偏离竖直方向的最大角度将小于 $45^{\circ}$。

\onepicture[align=right, w=4.0cm]{figs/24.png}
%% answer: 
\jdanswer{3 次}
%% memo:
\memoanswer{
设小球 $m$ 的摆线长度为 $l$。小球 $m$ 在下落过程中与 $M$ 相碰之前满足机械能守恒：

$mgl(1-\cos\theta)=\frac{1}{2}mv_{0}^{2}$ ①

$m$ 和 $M$ 碰撞过程满足：$mv_{0}=MV_{M}+mv_{1}$ ②

$\frac{1}{2}mv_{0}^{2}=\frac{1}{2}mv_{1}^{2}+\frac{1}{2}MV_{M}^{2}$ ③

联立②③得：$v_{1}=\frac{m-M}{m+M}v_{0}$ ④

说明小球被反弹，而后小球又以反弹速度和小球 $M$ 发生碰撞，满足：

$mv_{1}=MV_{M1}+mv_{2}$ ⑤

$\frac{1}{2}mv_{1}^{2}=\frac{1}{2}mv_{2}^{2}+\frac{1}{2}MV_{M1}^{2}$ ⑥

解得：$v_{2}=\frac{m-M}{m+M}|v_{1}|$ ⑦

整理得：$v_{2}=-\left(\frac{m-M}{m+M}\right)^{2}v_{0}$ ⑧

所以：$v_{n}=\left|\left(\frac{m-M}{m+M}\right)^{n}v_{0}\right|$ ⑨

而偏离方向为 $45^{\circ}$ 的临界速度满足：$mgl(1-\cos45^{\circ})=\frac{1}{2}mv_{\text{临界}}^{2}$ ⑩

联立①⑨⑩代入数据解得，当 $n=2$ 时，$v_{2}>v_{\text{临界}}$；当 $n=3$ 时，$v_{3}<v_{\text{临界}}$。所以，最多碰撞 3 次。
}

\item
%% number: 25
%% paperName: 2007年全国理综Ⅰ第 25 题 16 分
%% typeId: 6
%% type: 计算题
%% chapter: 磁场
%% point: 带电粒子在磁场中的运动
%% method: 无
%% score: 16
%% degree: 450
%% duplicateId: 0
%% body:
两屏幕荧光屏互相垂直放置，在两屏内分别取垂直于两屏交线的直线为 $x$ 和 $y$ 轴，交点 $O$ 为原点，如图所示。在 $y>0$，$0<x<a$ 的区域有垂直于纸面向内的匀强磁场，在 $y>0$，$x>a$ 的区域有垂直于纸面向外的匀强磁场，两区域内的磁感应强度大小均为 $B$。在 $O$ 点处有一小孔，一束质量为 $m$、带电量为 $q$（$q>0$）的粒子沿 $x$ 轴经小孔射入磁场，最后打在竖直和水平荧光屏上，使荧光屏发亮。入射粒子的速度可取从零到某一最大值之间的各种数值。已知速度最大的粒子在 $0<x<a$ 的区域中运动的时间与在 $x>a$ 的区域中运动的时间之比为 $2:5$，在磁场中运动的总时间为 $\frac{7}{12}T$，其中 $T$ 为该粒子在磁感应强度为 $B$ 的匀强磁场中做圆周运动的周期。试求两个荧光屏上亮线的范围（不计重力的影响）。

\onepicture[align=right, w=5.5cm]{figs/25.png}
%% answer: 
\jdanswer{$2a\leq x\leq 2\left(1+\frac{\sqrt{3}}{3}\right)a$}
%% memo:
\memoanswer{
对于 $y$ 轴上的光屏亮线范围的临界条件如图 1 所示：带电粒子的轨迹和 $x=a$ 相切，此时 $r=a$，$y$ 轴上的最高点为 $y=2r=2a$。

对于 $x$ 轴上光屏亮线范围的临界条件如图 2 所示：左边界的极限情况还是和 $x=a$ 相切，此刻，带电粒子在右边的轨迹是个圆，由几何知识得到在 $x$ 轴上的坐标为 $x=2a$；速度最大的粒子是如图 2 中的实线，由两段圆弧组成，圆心分别是 $c$ 和 $c'$，由对称性得到 $c'$ 在 $x$ 轴上，设在左右两部分磁场中运动时间分别为 $t_{1}$ 和 $t_{2}$，满足

$\frac{t_{1}}{t_{2}}=\frac{2}{5}$

$t_{1}+t_{2}=\frac{7}{12}T$

解得 $t_{1}=\frac{1}{6}T$，$t_{2}=\frac{5}{12}T$。由数学关系得到：$\sqrt{3}R=2a$，$OP=2a+R$，$OP=2\left(1+\frac{\sqrt{3}}{3}\right)a$。所以在 $x$ 轴上的范围是 $2a\leq x\leq 2\left(1+\frac{\sqrt{3}}{3}\right)a$。

\onepicture[w=5.0cm]{figs/25_an.jpg}
}

\end{enumerate}

\end{document}
